% generated by roam.py handbook export -- не редактировать, править заметки
% 20260922165110-функции_с_липшицевым_градиентом.org (body 581ec3dbcfaf)

\begin{definition}[{Функции с липшицевым градиентом}]\label[definition]{note:7DD1C7AB-1A49-4675-9F21-5A16C717F670}
  Это функции $f: E \rightarrow \R$, имеющие первые $L$-\hyperref[note:7E287F4B-AC3E-4CF5-9D32-3BC2BC52D01A]{гладкие производные} или $f \in C_L^{n,1}$ \\
  Можно записать так (\hyperref[note:1756907B-1FF2-4F62-8553-759B8509281C]{градиент}~--- $L$-липшицево отображение):
  \[
  \| \nabla f(x) - \nabla f(y)\|_2 \le L \| x - y\|_2\ \forall x,y \in E
  \]
\end{definition}

\begin{example}
  \begin{itemize}
    \item $f(x) = \frac{1}{2} x^T A x$, $A = A^T$: $\nabla f(x) = Ax$ и $\| \nabla f(x) - \nabla f(y) \| = \| A(x-y) \| \le \| A \|_2 \| x - y \|$, то есть $L = \| A \|_2$.
    \item $f(x) = x^2$: градиент $2x$ липшицев с $L = 2$, но сама $f$ не липшицева: $\abs{f(x) - f(0)} = x^2$ растет быстрее любой $L \abs{x}$. Поэтому «липшицева функция» и «функция с липшицевым градиентом»~--- разные понятия.
  \end{itemize}
\end{example}
